\section{Benchmark Design}

\begin{figure}[t]
\centering
\begin{tikzpicture}[font=\scriptsize, node distance=2.5mm and 2.5mm,
  box/.style={draw, rounded corners=2pt, align=center, text width=2.3cm, minimum height=1.3cm, inner sep=2.5pt, line width=0.6pt},
  wide/.style={box, text width=4.05cm, minimum height=1.0cm},
  arr/.style={->, line width=0.6pt, color=black!55}]
\node[box, draw=blue!45!black, fill=blue!4] (s1) {\textbf{1\; Scene family}\\[1pt] one task per sub-type: hazard, bystander, moving person};
\node[box, draw=blue!45!black, fill=blue!4, right=of s1] (s2) {\textbf{2\; Policy server}\\[1pt] GR00T N1.6, $\pi_{0.5}$, \ldots\ unmodified, remote; a new policy is a swap};
\node[box, draw=blue!45!black, fill=blue!4, right=of s2] (s3) {\textbf{3\; Recorders}\\[1pt] object pose, every link's pose (3-D body model), moving person};
\node[box, draw=red!55!black, fill=red!5, right=of s3] (s4) {\textbf{4\; Predicates T1--T6}\\[1pt] four dimensions, human-referenced; success-conditioned};
\node[box, draw=black!70, fill=black!5, right=of s4] (s5) {\textbf{5\; Report}\\[1pt] attempted / completing / violating, Wilson CI, exact tests};
\draw[arr] (s1) -- (s2); \draw[arr] (s2) -- (s3); \draw[arr] (s3) -- (s4); \draw[arr] (s4) -- (s5);
\node[wide, draw=orange!75!black, fill=orange!6, anchor=north west] (b1) at ([yshift=-5mm]s1.south west) {\textbf{6\; Fixability ablations}\\[1pt] name the hazard $\cdot$ hide it $\cdot$ safety command, at a placement where the rate can move};
\node[wide, draw=green!35!black, fill=green!5, right=of b1] (b2) {\textbf{7\; Feasibility witness}\\[1pt] a compliant completion exists in the scene --- an external layer as instrument, not a guard};
\node[wide, draw=black!60, fill=black!4, right=of b2] (b3) {\textbf{8\; Profile}\\[1pt] one unsafe rate per policy and sub-type (Table~\ref{tab:III}); new tasks plug in as scene families};
\draw[arr] (b1.north) -- (b1.north |- s1.south); \draw[arr] (b2.north) -- (b2.north |- s1.south); \draw[arr] (b3.north) -- (b3.north |- s1.south);
\end{tikzpicture}
\caption{\textbf{The benchmark protocol.} Scene family $\rightarrow$ unmodified remote policy $\rightarrow$ per-step recorders $\rightarrow$ per-sub-type predicates $\rightarrow$ success-conditioned report; fixability ablations and a feasibility witness make a cell a benchmark cell, and the cells form a policy $\times$ sub-type profile.}
\label{fig:pipeline}
\end{figure}




\subsection{Tasks, scenes and policies}


All GR00T tasks share one scene family: GR00T N1.6 \citep{ref5,ref13} drives a Unitree G1 in NVIDIA Isaac Sim via IsaacLab-Arena \citep{ref14,ref15} on a shelf-to-bin box carry --- pick a box from a shelf, walk $\approx$ 1.9 m down a corridor, release it in a bin --- and each sub-type changes what surrounds that carry (Fig. \ref{fig:gallery}): a hazard in the corridor --- a live strip, a hot stove or a standing person proxy (a 0.16 m-radius capsule plus a head sphere) --- for T1; a bystander beside the pick or bin zone at four positions for T2; a bystander at eight azimuths around the corridor for T3; the carried load's attitude for T4; a person standing on the path for T5a; and a person crossing the corridor --- a kinematic capsule with a collider and a contact sensor --- for T5b and T6. The static proxies carry no collider, so their ``contacts'' are geometric penetrations; only the crossing person can be struck. The corridor path is GR00T's own output: its navigation command is executed by a lower-body policy that sees neither language nor image (Appendix E.1). A second family puts the predicates around a Franka arm doing pick-and-place, driven by $\pi_{0.5}$ and $\pi_0$ (openpi), at a dining table, a kitchen counter and a packing station (Fig. \ref{fig:tabletop}): a keep-out between pick and place (T1); an adult (1.74 m) at the table, resting a forearm on it (T2, T5a); scissors and a fork, whose blade and tines give a real hazardous axis, with the person left or right (T3); a mug, with or without ``hot coffee, keep it upright'' (T4); a coworker's hand reaching into the destination bowl (T5b, T6). Per-step recorders log the payload's pose, every robot link's pose, the moving person's separation and the contact force.


\subsection{Metrics and attribution}


\textbf{Unsafe rate.} Each cell reports the fraction of episodes whose trajectory meets the sub-type's predicate, with a Wilson 95 \% interval. For transport sub-types the denominator is completing carries (box within 0.30 m of the bin): all 41 non-completing blind and hidden T1 episodes stall at the shelf, before the hazard is on the path (Fig. \ref{fig:t1uncond}), so conditioning removes no avoidance. T2 is scored over all episodes, since the sweep happens at the pick. Groups are compared with Fisher's exact test and paired conditions with McNemar's test; the episode is the unit throughout (Fig. \ref{fig:scatter} plots every cell's completion against its unsafe rate).

\textbf{Fixability ablations.} We name the hazard in the instruction, hide it from the cameras, or add an explicit safety command, on paired seeds. On the path these ablations face a ceiling --- the blind rate is already 100 \% --- so we also run them at a calibrated off-path placement where the blind rate is 37 \% (\S6).

\textbf{Attribution.} A rate is attributable to the policy only if the scene admits a compliant completion. We establish such \emph{feasibility witnesses} with external layers used as instruments, not proposed as guards: a repulsion shield given the hazard's coordinates (T1), a speed governor implementing the ISO/TS 15066 envelope together with that shield (T5a), and a protective stop (T5b, T6); their failure modes are reported with them (Appendix E). Table~\ref{tab:III} marks which rates have one; without it (T2, T3, the tabletop) the rate may be partly set by the scene. Appendix C gives thresholds and seeds, Appendix A every cell; completion under a substitute driver (2026-09-08 to 09-14) is not pooled (\S8).

